\section*{Capítulo \ref{ch:g12:oscillators-and-time} --- Mechanical Oscillators and the Measurement of Time}

\begin{solution}{exo:g12:oscillators-and-time:1}
(a) $T_0 = 2\pi\sqrt{1.00/9.81} = \qty{2.01}{s}$;
$f_0 = 1/T_0 = \qty{0.50}{Hz}$. (b) $T_0 \propto \sqrt{L}$: el
\omterm{def:g12:oscillators-and-time:oscillator}{período} dobles.
\end{solution}

\begin{solution}{exo:g12:oscillators-and-time:2}
Corazón: $f = 72/60 = \qty{1.2}{Hz}$, $T = \qty{0.83}{s}$. Ala:
$T = 1/30 = \qty{33}{ms}$. La elevación del pecho es la \emph{\omterm{def:g12:oscillators-and-time:oscillator}{amplitud}}.
\end{solution}

\begin{solution}{exo:g12:oscillators-and-time:3}
$T_0 = 2\pi\sqrt{0.10/25} = \qty{0.40}{s}$, $f_0 = \qty{2.5}{Hz}$.
Duplicación de $m$: $T_0 \times \sqrt{2} \approx 1.4$.
\end{solution}

\begin{solution}{exo:g12:oscillators-and-time:4}
\omterm{def:g12:oscillators-and-time:damping}{Pseudoperíodo} \qty{2.0}{s}. Relaciones $1.5/2.0 = 1.13/1.5 = 0.85/1.13
= 0.75$cada \omterm{def:g12:oscillators-and-time:oscillator}{período}; siguiente máximo$0.85 \times 0.75 \approx
\qty{0.64}{cm}$.
\end{solution}

\begin{solution}{exo:g12:oscillators-and-time:5}
Péndulo: $T_0 \propto 1/\sqrt{g}$, entonces $\times\sqrt{9.81/1.62} =
2.46$(más lento). Resorte--masa: sin cambios --- no$g$ en
$2\pi\sqrt{m/k}$.
\end{solution}

\begin{solution}{exo:g12:oscillators-and-time:6}
(a) $[L/g] = \unit{m}/(\unit{m/s^2}) = \unit{s^2}$; su raíz cuadrada
es un tiempo. (b) El \omterm{def:g10:orders-of-magnitude:unit}{kilogramo} aparece sólo en $m$: nada más podría
cancélelo, por lo que $m$ no puede ingresar. (c) Nunca: números puros como $2\pi$
son invisibles para las dimensiones.
\end{solution}

\begin{solution}{exo:g12:oscillators-and-time:7}
$\frac{dx}{dt} = -\frac{2\pi}{T_0}A\sin(2\pi t/T_0 + \varphi)$, entonces
$\frac{d^2x}{dt^2} = -(2\pi/T_0)^2 x(t)$. Igualdad con $-(k/m)x$
para todos $t$ se mantiene exactamente cuando $(2\pi/T_0)^2 = k/m$, es decir\
$T_0 = 2\pi\sqrt{m/k}$.
\end{solution}

\begin{solution}{exo:g12:oscillators-and-time:8}
$E_m = \tfrac12 kA^2 = \tfrac12 \times 40 \times 0.030^2 =
\qty{1.8e-2}{J}$;$v_{\max} = \sqrt{2E_m/m} = \sqrt{0.18} =
\qty{0.42}{m/s}$--- en$x = 0$; cero en los extremos$x = \pm A$.
\end{solution}

\begin{solution}{exo:g12:oscillators-and-time:9}
(a) $L = gT_0^2/4\pi^2 = \qty{0.994}{m}$. (b) $\Delta T_0/T_0 =
\tfrac12 \times 1.0\% = 0.5\%$:$86400 \times 0.005 \approx
\qty{4.3e2}{s}$ --- siete minutos lento por día.
\end{solution}

\begin{solution}{exo:g12:oscillators-and-time:10}
(a) $A \propto \sqrt{E}$: $\sqrt{0.95} \approx 0.975$, a $2.5\%$
caer. (b) $0.95^n = 0.5$: $n = \ln 2/\ln(1/0.95) \approx 14$
\omterm{def:g12:oscillators-and-time:oscillator}{periodos}. (c) No: para \omterm{def:g12:oscillators-and-time:damping}{mojadura} ligero se mantiene dentro de una fracción de
por ciento de $T_0$.
\end{solution}

\begin{solution}{exo:g12:oscillators-and-time:11}
(a) En $f_0$ cada impulso llega al paso y agrega \omterm{def:g10:energy-conservation:energy}{energía}; fuera de la cima,
empuja alternativamente ayuda y obstaculiza. (b) Marchar al paso es una
Unidad \omterm{def:g10:signals-and-waves:period}{periódico}; cerca del $f_0$ del puente subiría el \omterm{def:g12:oscillators-and-time:resonance}{resonancia}
pico. (c) El tambor barre el \omterm{def:g12:oscillators-and-time:resonance}{frecuencia natural} del marco.
durante el giro: momentáneo \omterm{def:g12:oscillators-and-time:resonance}{resonancia}, luego pasa el pico.
\end{solution}

\begin{solution}{exo:g12:oscillators-and-time:12}
(a) $T_0 \propto 1/\sqrt{g}$: lento, por $\Delta T_0/T_0 = \tfrac12
\times 0.03/9.81 = \num{1.5e-3}$, es decir $\approx \qty{132}{s}$ por
día. (b) $L \propto g$ en $T_0$ fijo: acortar en $0.994 \times
0.03/9.81 \approx \qty{3.0}{mm}$.
\end{solution}

\begin{solution}{exo:g12:oscillators-and-time:13}
$v_{\max} = 2\pi A/T_0 = 2\pi \times 0.040/0.63 = \qty{0.40}{m/s}$;
$a_{\max} = (2\pi/T_0)^2 A = \qty{4.0}{m/s^2}$; $k = m(2\pi/T_0)^2
= \qty{20}{N/m}$;$E_m = \tfrac12 kA^2 = \qty{1.6e-2}{J}$.
\end{solution}

\begin{solution}{exo:g12:oscillators-and-time:14}
$T_0 = 40.1/20 = \qty{2.005}{s}$; $g = 4\pi^2 L/T_0^2 = 4\pi^2
\times 0.995/4.020 = \qty{9.77}{m/s^2}$. Incertidumbre$\approx 0.2\%
+ 2 \times 0.5\% = 1.2\%$, es decir $\pm\qty{0.12}{m/s^2}$: el
El intervalo contiene \qty{9.81}{m/s^2} --- consistente.
\end{solution}

\begin{solution}{exo:g12:oscillators-and-time:15}
Péndulo $10/86400 \approx \num{1.2e-4}$; cuarzo $0.5/\num{2.6e6}
\approx \num{1.9e-7}$; cesio$1/\num{3.2e15} \approx \num{3e-16}$.
Un error \qty{1}{\micro s} es $\num{1000} \times \qty{30}{cm} =
\qty{300}{m}$ de posición: el GPS necesita la familia atómica.
\end{solution}

\begin{solution}{pb:g12:oscillators-and-time:1}
\textbf{1.} El error de reacción $\approx \qty{0.2}{s}$ se divide por
$50$: sobre \qty{4}{ms} en $T_0$.

\textbf{2.} $T = 99.4/50 = \qty{1.988}{s}$.

\textbf{3.} $L = gT^2/4\pi^2 = 9.81 \times 1.988^2/4\pi^2 \approx
\qty{0.982}{m}$.

\textbf{4.} $L = \qty{0.994}{m}$ (el péndulo de segundos): bajar el
bob, alargándose en $\approx \qty{1.2}{cm}$.

\textbf{5.} Demasiado corto, demasiado rápido: rápido, por $(2.000 -
1.988)/1.988 \times 86400 \approx \qty{5.2e2}{s} \approx 8.7$
minutos por día.

\textbf{6.} Luna: $\times\sqrt{9.81/1.62} = 2.46$ --- inutilizable.
Observatorio: $\Delta T_0/T_0 = \tfrac12 \times 0.01/9.81 =
\num{5.1e-4}$: pierde$\approx \qty{44}{s}$ por día.

\textbf{7.} $\Delta L = L\lambda\,\Delta\theta = 0.994 \times
\num{1.9e-5} \times 5.0 \approx \qty{9.4e-5}{m} \approx
\qty{0.1}{mm}$.

\textbf{8.} $T' = 2\pi\sqrt{L(1 + \lambda\Delta\theta)/g} =
T_0\sqrt{1 + \lambda\Delta\theta} \approx T_0\,(1 + \tfrac12
\lambda\Delta\theta)$.

\textbf{9.} $\Delta T_0/T_0 = \tfrac12 \times \num{1.9e-5} \times
5.0 = \num{4.8e-5}$:$86400 \times \num{4.8e-5} \approx
\qty{4.1}{s}$ por día, lento (varilla más larga, \omterm{def:g12:oscillators-and-time:oscillator}{período} más larga).

\textbf{10.} $\Delta\theta = \qty{-10}{K}$: $\Delta T_0/T_0 =
\num{-9.5e-5}$, gana$\approx \qty{8.2}{s}$ por día.

\textbf{11.} El acero y el latón se expanden de manera diferente; colgado en alternancia
direcciones, las varillas de latón bajan la pesa mientras que las de acero suben
eso. Dimensionado para que los dos turnos se cancelen, la longitud efectiva --- y el
tasa --- sobrevive a las estaciones.

\textbf{12.} $280/2.000 = 140$ \omterm{def:g12:oscillators-and-time:oscillator}{periodos}.

\textbf{13.} $\theta_m = 2.00^\circ = \qty{0.0349}{rad}$: $E_m =
\tfrac12 \times 1.2 \times 9.81 \times 0.994 \times 0.0349^2 \approx
\qty{7.1e-3}{J}$.

\textbf{14.} Por \omterm{def:g12:oscillators-and-time:oscillator}{período} el \omterm{def:g12:oscillators-and-time:oscillator}{amplitud} se reduce en $2^{-1/140} \approx
0.9951$, el \omterm{def:g10:energy-conservation:energy}{energía} en$0.9951^2 \approx 0.990$: alrededor de$1.0\%$ de
$E_m$ por \omterm{def:g12:oscillators-and-time:oscillator}{período}.

\textbf{15.} Empuje $\approx 0.010 \times \num{7.1e-3} =
\qty{7.1e-5}{J}$; \omterm{def:g11:work-of-force:power}{fuerza}$\num{7.1e-5}/2.0 \approx \qty{3.5e-5}{W}$
--- decenas de microvatios.

\textbf{16.} Una semana es $604800/2.0 = \num{3.02e5}$ \omterm{def:g12:oscillators-and-time:oscillator}{periodos}, necesitando
$\num{3.02e5} \times \num{7.1e-5} \approx \qty{21}{J}$; el \omterm{def:g10:universal-gravitation:weight}{peso}
suministros $Mgh = 2.5 \times 9.81 \times 1.0 = \qty{24.5}{J}$: suficiente,
en $21/24.5 \approx 88\%$ \omterm{def:g10:energy-conservation:efficiency}{eficiencia}.

\textbf{17.}$T = 1/\num{32768} = \qty{30.5}{\micro s}$. $\num{32768}
= 2^{15}$: quince etapas de división por dos hacen girar el cristal \omterm{def:g10:signals-and-waves:signal}{señal}
en un tick exacto \qty{1}{Hz}.

\textbf{18.} Cuarzo: $0.3/\num{2.6e6} \approx \num{1.2e-7}$.
Péndulo: $0.5/86400 \approx \num{5.8e-6}$, es decir $\approx
\qty{15}{s}$por mes. Quartz gana por un factor$\approx 50$.

\textbf{19.} $1/\num{3.2e15} \approx \num{3e-16}$: nueve órdenes de
magnitud por debajo del cuarzo. Tiempos de GPS \omterm{def:g10:signals-and-waves:signal}{señales} a nanosegundos ---
\qty{30}{cm} de luz viajan cada uno --- por lo que solo los relojes atómicos servirán.

\textbf{20.} Restaurado, el reloj de pared tiene capacidad para $\pm\qty{15}{s}$ por
mes contra el $\pm\qty{0.3}{s}$ del cuarzo: guárdelo para el timbre,
Confía en el cuarzo para el tren --- y el segundo ahora pertenece
al cesio \omterm{def:g11:fundamental-interactions:ladder}{átomo}.
\end{solution}